% This document was generated by an LLM.

\documentclass{workbook}

\DeclareMathOperator{\arcsinh}{arcsinh}
\DeclareMathOperator{\arccosh}{arccosh}
\DeclareMathOperator{\arctanh}{arctanh}
\DeclareMathOperator{\sech}{sech}
\DeclareMathOperator{\gd}{gd}
\newcommand{\R}{\mathbb{R}}
\newcommand{\C}{\mathbb{C}}

\title{Circular and Hyperbolic Functions\\
\large A Connections and Proof Workbook}
\author{}
\date{}

\begin{document}
\maketitle

\noindent
This workbook builds the connection between circular and hyperbolic functions
through geometry, differential equations, exponentials, inverse functions,
and calculus. It contains questions and proof requests, not solutions. The
problems are ordered so that later sections can use results proved earlier.
The complex-variable sections are approachable with elementary complex
arithmetic and power series; a course in complex analysis is not assumed.

\medskip\noindent
\textbf{Working convention.} Angles are in radians. Unless stated otherwise,
variables are real. Write $\arcsinh$, $\arccosh$, and $\arctanh$ for inverse
hyperbolic functions. You may use the standard definitions and derivatives of
real $\sin$, $\cos$, $e^x$, and $\ln x$, but derive hyperbolic identities when
asked. For complex problems, use $i^2=-1$ and the power-series definition
$e^z=\sum_{n=0}^{\infty}z^n/n!$; you may use $e^{z+w}=e^ze^w$ and
$e^{it}=\cos t+i\sin t$. A complex square root or logarithm is multivalued
unless a branch is specified. When a problem says ``principal,'' take
$\operatorname{Arg}z\in(-\pi,\pi)$ on the plane cut along the nonpositive
real axis, $\operatorname{Log}z=\ln|z|+i\operatorname{Arg}z$, and
$\sqrt z=\exp(\tfrac12\operatorname{Log}z)$ there. Boundary values on a cut
must be stated explicitly.

\medskip\noindent
\textbf{How to work.} For a recall problem, state domains, ranges, and branch
restrictions. For a proof, name the identity or earlier result you use. For a
``decide'' problem, prove the statement or give a counterexample. In problems
with complex inverses, first say where the formula is meant to hold.

\section{The two real families}

\begin{problem}{Recall.}
State the definitions of $\sinh t$, $\cosh t$, and $\tanh t$ in terms of
$e^t$. What are their real domains and ranges?
\end{problem}

\begin{problem}{Prove.}
Use the exponential definitions to establish the parity of $\sinh$ and
$\cosh$, then derive $\cosh^2t-\sinh^2t=1$ for every $t\in\R$.
\end{problem}

\begin{problem}{Compare.}
Place $\cos^2t+\sin^2t=1$ beside the hyperbolic identity. Explain what
each equation says about the point $(x(t),y(t))$ on its corresponding curve.
Which point of each curve occurs at $t=0$?
\end{problem}

\begin{problem}{Prove.}
Derive the real derivatives of $\sinh$, $\cosh$, and $\tanh$ from their
definitions. In the derivative of $\tanh$, identify exactly where the
hyperbolic Pythagorean identity enters.
\end{problem}

\begin{problem}{Prove.}
Determine the intervals of increase and decrease of $\sinh$, $\cosh$,
and $\tanh$. Find their limits as $t\to\pm\infty$ and justify the ranges
stated in Problem~1.
\end{problem}

\begin{problem}{Decide.}
Is $\cosh t\geq|\sinh t|$ true for all real $t$? Can equality occur?
Give a proof or counterexample to each claim.
\end{problem}

\begin{problem}{Prove.}
Show that $\cosh t\geq1$ for all real $t$, with equality exactly when
$t=0$. Give one proof using the hyperbolic identity and another using
elementary facts about $e^t$.
\end{problem}

\begin{problem}{Decide.}
For each assertion, determine whether it holds for all real $t$:
$\sinh t\leq\cosh t$, $|\tanh t|<1$, and $\cosh t=\cos t$.
For a false assertion, find the exact set of $t$ for which it holds.
\end{problem}

\section{Addition laws and differential equations}

\begin{problem}{Prove.}
Derive formulas for $\sinh(s+t)$ and $\cosh(s+t)$ directly from
$e^{s+t}=e^se^t$. Compare the sign in the $\cosh$ formula with the sign
in the circular cosine addition formula.
\end{problem}

\begin{problem}{Prove.}
Derive the double-angle formulas for $\sinh(2t)$ and $\cosh(2t)$.
Write $\cosh(2t)$ in two further ways using
$\cosh^2t-\sinh^2t=1$.
\end{problem}

\begin{problem}{Prove.}
Find an addition formula for $\tanh(s+t)$, giving the conditions under
which each quotient is defined. Show that its denominator is positive
when $s,t\in\R$.
\end{problem}

\begin{problem}{Recall and prove.}
Which second-order differential equation is satisfied by each of
$\sin t$, $\cos t$, $\sinh t$, and $\cosh t$? Verify the four answers by
differentiation, and record the initial values and initial derivatives.
\end{problem}

\begin{problem}{Investigate.}
Let $J=\left(\begin{smallmatrix}0&-1\\1&0\end{smallmatrix}\right)$ and
$K=\left(\begin{smallmatrix}0&1\\1&0\end{smallmatrix}\right)$.
Compute $J^2$ and $K^2$. How do these two signs predict the different
second-order equations in the preceding problem?
\end{problem}

\begin{problem}{Prove.}
Define
$R(t)=\left(\begin{smallmatrix}\cos t&-\sin t\\\sin t&\cos t\end{smallmatrix}\right)$
and
$B(t)=\left(\begin{smallmatrix}\cosh t&\sinh t\\\sinh t&\cosh t\end{smallmatrix}\right)$.
Show that $R(s)R(t)=R(s+t)$ and $B(s)B(t)=B(s+t)$.
What are $R(0)$, $B(0)$, and their inverses?
\end{problem}

\begin{problem}{Prove.}
For a vector $(x,y)\in\R^2$, show that $R(t)$ preserves $x^2+y^2$
and that $B(t)$ preserves $x^2-y^2$. Calculate both determinants.
Explain why determinant $1$ alone would not prove either preservation law.
\end{problem}

\begin{problem}{Synthesis.}
Starting with the differential equations and the initial values, explain
why one coordinate pair travels around a circle while the other travels
along a hyperbola. Identify which earlier identities your explanation uses.
\end{problem}

\section{Area and the meaning of angle}

\medskip\noindent
For this section you may use the signed swept-area formula
\[
  A(t)=\frac12\int_0^t
  \bigl(x(s)y'(s)-y(s)x'(s)\bigr)\,ds
\]
for a differentiable curve $(x(s),y(s))$ beginning at $(1,0)$.

\begin{problem}{Prove.}
Apply the swept-area formula to
$(x(t),y(t))=(\cos t,\sin t)$. Determine $A(t)$ for positive and
negative $t$. What changes if area is required to be unsigned?
\end{problem}

\begin{problem}{Prove.}
Apply the same formula to
$(x(t),y(t))=(\cosh t,\sinh t)$. Identify exactly which identity
makes the integrand constant.
\end{problem}

\begin{problem}{Compare.}
Compute the speed of the unit-circle and unit-hyperbola
parametrizations. Is $|t|$ the Euclidean arc length from parameter $0$
to parameter $t$ for either curve? Justify both answers.
\end{problem}

\begin{problem}{Concept.}
Explain why the parameter $t$ is called a hyperbolic angle even though
it is not Euclidean arc length. What feature makes real circular angle
periodic, and what feature makes real hyperbolic angle unbounded?
\end{problem}

\begin{problem}{Decide.}
Does the substitution $y\mapsto iy$ turn $x^2+y^2=1$ into
$x^2-y^2=1$ as an algebraic equation? Does it describe an ordinary
rotation in the real $(x,y)$-plane? Explain the distinction.
\end{problem}

\section{One complex structure}

\begin{problem}{Recall.}
State the exponential formulas for $\sin z$ and $\cos z$ for
$z\in\C$. Then state the exponential formulas for $\sinh z$ and
$\cosh z$. Specify which of these functions are defined for every
complex number.
\end{problem}

\begin{problem}{Prove.}
Starting from the exponential formulas, derive
$\cos(iz)=\cosh z$ and $\sin(iz)=i\sinh z$ for $z\in\C$.
Reverse the substitution to obtain two converse formulas.
\end{problem}

\begin{problem}{Prove.}
Where both sides are defined, obtain the corresponding identities for
$\tan(iz)$ and $\tanh(iz)$. State explicitly why the phrase ``where
both sides are defined'' is needed for these quotients.
\end{problem}

\begin{problem}{Prove.}
Use the power series of $\sin z$, $\cos z$, $\sinh z$, and $\cosh z$
to rederive the two rotation identities proved above. Identify the role of
$i^2=-1$ in the odd and even terms separately.
\end{problem}

\begin{problem}{Translate.}
Starting with $\cos^2z+\sin^2z=1$, substitute $iz$ for $z$ and
simplify. Repeat with the circular formulas for $\cos(2z)$ and
$\sin(2z)$. Compare with your results from Section~2.
\end{problem}

\begin{problem}{Prove.}
Using the addition formulas, express $\sin(x+iy)$ and $\cos(x+iy)$
in terms of real circular functions of $x$ and real hyperbolic functions
of $y$.
\end{problem}

\begin{problem}{Prove.}
Find $|\sin(x+iy)|^2$ and $|\cos(x+iy)|^2$ for real $x,y$.
Use either formula to explain why boundedness of $\sin x$ and $\cos x$
on $\R$ does not imply boundedness on $\C$.
\end{problem}

\begin{problem}{Compute.}
Evaluate $\sin(i\ln2)$, $\cos(i\ln2)$,
$\sin(\pi/2+iy)$, and $\cos(\pi/2+iy)$ for real $y$.
Keep the answers exact.
\end{problem}

\begin{problem}{Decide.}
Does $\sin(iz)=i\sinh z$ mean that $\sin x=\sinh x$ for real $x$?
Does it mean that the hyperbolic functions are restrictions of the
complex circular functions to the imaginary axis? State the precise
version of the second claim, including factors of $i$.
\end{problem}

\section{Real inverse functions}

\begin{problem}{Recall.}
State the restricted real domains used to define $\arcsin$, $\arccos$,
and $\arctan$. Give their resulting ranges and the domains of the
inverse functions.
\end{problem}

\begin{problem}{Recall.}
State the real domains and ranges of $\arcsinh$, $\arccosh$, and
$\arctanh$. Which original hyperbolic function must be restricted
before it can be inverted, and to what interval?
\end{problem}

\begin{problem}{Prove.}
Let $u=e^w>0$. Starting from $\sinh w=x$, solve the resulting
quadratic for $u$ and derive a real logarithmic formula for
$\arcsinh x$. Justify the choice of root for every $x\in\R$.
\end{problem}

\begin{problem}{Prove.}
Use the same substitution to derive the real logarithmic formula for
$\arccosh x$ when $x\geq1$. Explain how the two quadratic roots
correspond to the evenness of $\cosh$.
\end{problem}

\begin{problem}{Prove.}
Derive the logarithmic formula for $\arctanh x$ when $|x|<1$.
Why is this domain necessary for a real inverse? Check the signs of
every quantity inside the logarithm.
\end{problem}

\begin{problem}{Compute.}
Find exact values of $\arcsinh(3/4)$, $\arccosh(5/4)$,
and $\arctanh(3/5)$. Verify each answer by composing with the
corresponding hyperbolic function.
\end{problem}

\begin{problem}{Decide.}
Evaluate $\arcsin(\sin(5\pi/6))$ and
$\arccosh(\cosh(-2))$. For each, explain why the inverse does not
return the original input.
\end{problem}

\begin{problem}{Prove.}
For real $x$, show $\arcsinh(-x)=-\arcsinh x$.
Prove the analogous claim for $\arctanh$ on its real domain.
Can the same claim hold for $\arccosh$? Give a domain-based reason.
\end{problem}

\section{Complex inverses and branch choices}

\begin{problem}{Investigate.}
Solve $e^w=1$ and $e^w=-1$ for all $w\in\C$. Explain from these
solutions why the complex logarithm has infinitely many values.
\end{problem}

\begin{problem}{Derive.}
Set $u=e^{iw}$ and solve $\sin w=z$ as a quadratic equation in $u$.
Write the resulting candidates for $w$ using a logarithm, and identify
both kinds of choices involved before an inverse is single-valued.
\end{problem}

\begin{problem}{Derive.}
Repeat the preceding problem for $\cos w=z$. For real
$x\in[-1,1]$, explain how to select $w\in[0,\pi]$.
\end{problem}

\begin{problem}{Derive.}
Solve $\tan w=z$ using $u=e^{iw}$. On the real axis, verify that
your selected value lies in $(-\pi/2,\pi/2)$ when $z\in\R$.
\end{problem}

\begin{problem}{Prove with a domain.}
Starting from $\sin(iw)=i\sinh w$, derive a relation between
$\arcsin z$ and $\arcsinh(iz)$ near $z=0$. State why an unqualified
global assertion about principal inverses would require checking
branch cuts.
\end{problem}

\begin{problem}{Prove with a domain.}
Derive the analogous local relation between $\arctan z$ and
$\arctanh(iz)$. Identify the branch points of both inverse
functions from their logarithmic formulas.
\end{problem}

\begin{problem}{Audit.}
A calculation replaces $\operatorname{Log}(ab)$ by
$\operatorname{Log}a+\operatorname{Log}b$ for arbitrary nonzero complex
$a,b$. Find values away from the logarithm's cut for which this is
false, and explain the discrepancy.
\end{problem}

\begin{problem}{Audit.}
Someone applies the real formula
$\arccosh x=\ln(x+\sqrt{x^2-1})$ at the complex input $z=-2$
using the positive square root of $z^2-1$ and a boundary value of the
logarithm. Compare the resulting real part with that of the principal
$\arccosh$ boundary value. Where did the
extension of the real formula become unreliable?
\end{problem}

\begin{problem}{Concept.}
What is a branch cut intended to accomplish? Explain why the
square root and the logarithm each need a branch choice in formulas
for inverse circular and inverse hyperbolic functions.
\end{problem}

\section{Derivatives, integrals, and substitutions}

\begin{problem}{Prove.}
Differentiate $\arcsinh x$, $\arccosh x$, and $\arctanh x$ by
implicit differentiation. State each real domain on which the
derivative formula holds, and justify every square-root sign.
\end{problem}

\begin{problem}{Compare.}
Derive the real derivatives of $\arcsin x$ and $\arctan x$ and set
them beside the hyperbolic inverse derivatives. Describe the exact
sign changes in the radicands or denominators.
\end{problem}

\begin{problem}{Integrate.}
Find real antiderivatives of $1/(1+x^2)$ on $\R$ and of
$1/(1-x^2)$ on each interval avoiding $\pm1$. Explain precisely
where the latter antiderivative can be named $\arctanh x$.
\end{problem}

\begin{problem}{Integrate.}
For $a>0$, evaluate
$\int dx/\sqrt{a^2+x^2}$ by a hyperbolic substitution.
Give the answer in both inverse-function and logarithmic form.
\end{problem}

\begin{problem}{Compare substitutions.}
For $a>0$, simplify $\sqrt{a^2+x^2}$ first with $x=a\sinh t$
and then with $x=a\tan\theta$. State real restrictions needed to
remove absolute values from the square roots. What is the advantage
of each substitution in the preceding integral?
\end{problem}

\begin{problem}{Integrate.}
For $a>0$ and $x>a$, find an antiderivative of
$1/\sqrt{x^2-a^2}$ using $x=a\cosh t$. Explain why the same
formula cannot be applied unchanged to $x<-a$ without checking a sign.
\end{problem}

\begin{problem}{Decide.}
Is $\sqrt{x^2}=x$ valid for every real $x$? Locate an inverse-function
derivative or substitution in this section for which that error would
lead to an incorrect sign.
\end{problem}

\section{The Gudermannian and another real bridge}

\medskip\noindent
For this section define $\gd t=\arctan(\sinh t)$ for $t\in\R$.

\begin{problem}{Prove.}
Determine the range and parity of $\gd$. Differentiate it and express
its derivative using $\sech t=1/\cosh t$.
\end{problem}

\begin{problem}{Prove.}
Derive formulas for $\tan(\gd t)$, $\sec(\gd t)$, and
$\sin(\gd t)$ in terms of hyperbolic functions. State the sign
argument that selects the correct value of the secant.
\end{problem}

\begin{problem}{Prove.}
Show that $\gd t=\int_0^t\sech s\,ds$.
Why is the constant of integration fixed by this equation?
\end{problem}

\begin{problem}{Prove.}
Find $\gd^{-1}\theta$ for $|\theta|<\pi/2$, first as an inverse
hyperbolic function and then as a real logarithm. Differentiate
your formula to obtain an antiderivative of $\sec\theta$.
\end{problem}

\begin{problem}{Investigate.}
The vertical coordinate of the spherical Mercator map at latitude
$\theta$ is proportional to $\int_0^\theta\sec u\,du$.
Explain the role of $\gd^{-1}$ in this formula and determine what
happens as $\theta\to\pm\pi/2$.
\end{problem}

\begin{problem}{Synthesis.}
Compare the two bridges $\sin(it)=i\sinh t$ and
$\sin(\gd t)=\tanh t$. What changes in the input in each case?
Which relation stays entirely within real variables?
\end{problem}

\section{Proof audits and synthesis}

\begin{problem}{Audit.}
``$\cosh^2t-\sinh^2t=1$, so $\cosh t=\sqrt{1+\sinh^2t}$ for
all complex $t$.'' Identify the step that needs a branch or sign
choice. State a correct real-variable version.
\end{problem}

\begin{problem}{Audit.}
``The circle and the hyperbola are both parametrized by arc length
$t$ because each swept area is $t/2$.'' Locate the invalid inference
and calculate a quantity that settles it.
\end{problem}

\begin{problem}{Audit.}
``$\arccosh(\cosh t)=t$ for every real $t$ because inverse
functions undo each other.'' State the exact result that is true
and justify the necessary restriction.
\end{problem}

\begin{problem}{Audit.}
``Substitute $it$ into $1+\tan^2t=\sec^2t$ and replace the names
to get $1+\tanh^2t=\sech^2t$.'' Identify the missing factor and
derive the correct identity.
\end{problem}

\begin{problem}{Decide.}
For real $x$, are $\sin x$ and $\sinh x$ ever equal apart from
$x=0$? Provide a proof or a counterexample to the claim that
$x=0$ is their only intersection.
\end{problem}

\begin{problem}{Synthesis.}
Give two independent derivations of
$\cosh^2t-\sinh^2t=1$: one using exponentials and one using a
differential equation or preserved quadratic form. Explain what
each derivation reveals about the circular connection.
\end{problem}

\begin{problem}{Synthesis.}
Choose one circular and one hyperbolic inverse. Derive their
logarithmic forms from exponentials, state their real domains,
and explain why the corresponding complex formulas need branch
choices. Make every choice explicit.
\end{problem}

\begin{problem}{Synthesis.}
Write a short explanation for a student who has not studied
complex analysis: in what precise sense are circular and
hyperbolic functions connected, and which part of the connection
can be understood using only real calculus? Include one formula,
one geometric fact, and one limitation of the analogy.
\end{problem}

\section*{Final self-check}

Without consulting notes, reconstruct the two real parametrizations and
their preserved quadratic forms. Derive the $z\mapsto iz$ identities from
exponentials, explain one inverse logarithm from its quadratic equation,
and name the branch choice that prevents a naive global inversion.
Finally, give a real-variable explanation of the Gudermannian and identify
one integration problem where the hyperbolic form is useful.

\end{document}
